\chapter{Information and Relational Order}

The temporal programme begins from relational states rather than from a pre-existing time coordinate. Each state $s$ is assigned a probability distribution
\[
p(s)=(p_1,\ldots,p_m),
\qquad
p_i\geq0,
\qquad
\sum_i p_i=1,
\]
and its Shannon entropy in bits,
\[
\boxed{
H_S(s)=-\sum_i p_i(s)\log_2p_i(s).
}
\]
The inherited normalization is
\[
\kappa=\frac{\ln2}{24\pi}.
\]

For an admitted directed transition $e:a\to b$, the first informational quantity is the exact state-function difference
\[
\Delta H_e=H_S(b)-H_S(a).
\]
This quantity measures the change in the Shannon state entropy across that transition. Importantly, it is an exact difference. Therefore, for every closed directed cycle $C$,
\[
\boxed{
\sum_{e\in C}\Delta H_e=0.
}
\]
The identity follows by telescoping and has a direct methodological consequence: a state entropy difference can contribute to an open transition, but it cannot by itself generate a non-zero circulation around a closed cycle.

\section{Complex relational state and phase transport}

The same relational state is also represented by a normalized complex vector
\[
\psi(s)\in\mathbb C^m,
\qquad
\langle\psi|\psi\rangle=1.
\]
Whenever two consecutive states have non-zero overlap, define the normalized Pancharatnam link
\[
G_{a\to b}
=
\frac{\langle\psi_a|\psi_b\rangle}
{|\langle\psi_a|\psi_b\rangle|}
\in U(1).
\]
An open link is gauge-covariant under local rephasing of the state representatives. The ordered product of links around a closed cycle is gauge-invariant and defines the discrete geometric holonomy
\[
\gamma_B(C)
=
\operatorname{Arg}\!\left(\prod_{e\in C}G_e\right).
\]

\section{Shannon--phase transition link}

To connect informational change and phase transport without identifying them prematurely, introduce a transition-associated non-exact information increment $\sigma_e$. The first microscopic candidate closure is obtained from forward/reverse transition asymmetry. For strictly positive admitted transition probabilities,
\[
\boxed{
\sigma_{a\to b}
=
\log_2\frac{P(b\mid a)}{P(a\mid b)}.
}
\]
This quantity is path-level rather than a state entropy. It obeys $\sigma_{b\to a}=-\sigma_{a\to b}$ and vanishes for a symmetric transition pair. The candidate temporal link is
\[
\boxed{
L_e
=
G_e\exp\!\left[i\kappa\left(\Delta H_e+\sigma_e\right)\right].
}
\]
The corresponding transition phase weight is
\[
q_e=\operatorname{Arg}L_e.
\]
It can therefore enter the atomic temporal transition measure used later in the construction,
\[
\mathcal K_T=\sum_e q_e\,\delta_{s_e}.
\]

For a closed cycle the exact Shannon contribution cancels. Defining the cycle affinity
\[
\mathcal A_C
=
\sum_{e\in C}\sigma_e
=
\log_2\prod_{e\in C}
\frac{P_{e,+}}{P_{e,-}},
\]
gives
\[
\boxed{
\operatorname{Arg}\!\left(\prod_{e\in C}L_e\right)
=
\gamma_B(C)
+
\kappa\mathcal A_C
\pmod{2\pi}.
}
\]
This is the first structural separation in the derivation. The geometric part supplies phase holonomy. The exact Shannon state potential supplies open-transition information change. Directed path affinity supplies a non-exact information channel that can survive cycle closure.

A useful control follows immediately. If a strictly positive relational weight $\pi_s$ satisfies pairwise detailed balance,
\[
\pi_aP(b\mid a)=\pi_bP(a\mid b),
\]
then
\[
\sigma_{a\to b}=\log_2\frac{\pi_b}{\pi_a}
\]
and the cycle affinity telescopes to zero. Non-zero $\mathcal A_C$ is therefore a precise obstruction to this detailed-balance representation on the declared cycle.

The next derivational target is to determine how the relational density and viscosity fields modify the forward/reverse transition weights, rather than assigning those weights by hand.
